\documentclass[12pt]{article}
\usepackage[utf8]{inputenc}
\usepackage[spanish]{babel}
\usepackage{amsmath, amssymb}
\usepackage{geometry}
\geometry{a4paper, margin=1in}

\title{Cinemática de la Rotación y la Matriz Antisimétrica en el Efecto Dzhanibekov}
\author{Hugo Sosa}
\date{\today}

\begin{document}

\maketitle

\section{Introducción}
Para simular el efecto Dzhanibekov, resolvemos las ecuaciones de Euler que nos dan las componentes de la velocidad angular en el sistema de referencia del cuerpo ($\omega_1, \omega_2, \omega_3$). Sin embargo, para visualizar el cuerpo rotando en el espacio real (sistema de laboratorio o fijo), necesitamos una forma de actualizar su orientación.

En el código, utilizamos una matriz de rotación $R$ y su derivada temporal $\dot{R}$. La clave de la implementación reside en la relación:
\begin{equation}
    \dot{R} = R \Omega_{body}
\end{equation}
donde $\Omega_{body}$ es una matriz antisimétrica construida a partir de $\vec{\omega}$. En este documento explicaremos de dónde sale esta matriz y por qué funciona.

\section{La Matriz de Rotación como un Conjunto de Ejes}
Imagina que el sistema de laboratorio tiene ejes fijos $\hat{x}, \hat{y}, \hat{z}$. El cuerpo rígido tiene sus propios ejes principales de inercia $\hat{e}_1, \hat{e}_2, \hat{e}_3$. La matriz de rotación $R$ en cualquier instante es simplemente una matriz cuyas columnas son estos ejes principales expresados en las coordenadas del laboratorio:
\begin{equation}
    R = \begin{bmatrix} | & | & | \\ \hat{e}_1 & \hat{e}_2 & \hat{e}_3 \\ | & | & | \end{bmatrix}
\end{equation}

\section{La Dinámica de los Ejes}
De la mecánica clásica básica, sabemos que si un vector $\vec{A}$ está "pegado" a un cuerpo que rota con velocidad angular $\vec{\omega}$, su derivada temporal vista desde el sistema fijo es:
\begin{equation}
    \frac{d\vec{A}}{dt} = \vec{\omega}_{fijo} \times \vec{A}
\end{equation}
Como nuestros ejes $\hat{e}_i$ están pegados al cuerpo, su evolución es:
\begin{equation}
    \dot{\hat{e}}_1 = \vec{\omega}_{fijo} \times \hat{e}_1, \quad \dot{\hat{e}}_2 = \vec{\omega}_{fijo} \times \hat{e}_2, \quad \dot{\hat{e}}_3 = \vec{\omega}_{fijo} \times \hat{e}_3
\end{equation}
Esto significa que la derivada de la matriz $R$ es:
\begin{equation}
    \dot{R} = \begin{bmatrix} | & | & | \\ \vec{\omega}_{fijo} \times \hat{e}_1 & \vec{\omega}_{fijo} \times \hat{e}_2 & \vec{\omega}_{fijo} \times \hat{e}_3 \\ | & | & | \end{bmatrix}
\end{equation}

\section{El Producto Vectorial como Matriz}
Aquí es donde entra la magia del álgebra lineal. El producto vectorial $\vec{a} \times \vec{b}$ se puede escribir como el producto de una matriz antisimétrica por un vector:
\begin{equation}
    \vec{a} \times \vec{b} = [\vec{a}]_\times \vec{b} = \begin{bmatrix} 0 & -a_z & a_y \\ a_z & 0 & -a_x \\ -a_y & a_x & 0 \end{bmatrix} \begin{bmatrix} b_x \\ b_y \\ b_z \end{bmatrix}
\end{equation}
A esta matriz $[\vec{a}]_\times$ se la llama a veces matriz \textit{skew-symmetric} o matriz antisimétrica.

\section{Cuerpo vs. Laboratorio}
En el código, tenemos $\vec{\omega}_{cuerpo} = (\omega_1, \omega_2, \omega_3)$. La relación entre la velocidad angular en el sistema fijo y en el sistema del cuerpo es:
\begin{equation}
    \vec{\omega}_{fijo} = R \vec{\omega}_{cuerpo}
\end{equation}

Si usamos esta relación y la propiedad del producto vectorial bajo rotaciones, se puede demostrar (tras un poco de álgebra) que la evolución de $R$ se puede expresar de dos formas equivalentes:
\begin{enumerate}
    \item Usando $\vec{\omega}_{fijo}$: $\dot{R} = [\vec{\omega}_{fijo}]_\times R$
    \item Usando $\vec{\omega}_{cuerpo}$: $\dot{R} = R [\vec{\omega}_{cuerpo}]_\times$
\end{enumerate}

\section{¿Por qué usamos la segunda opción?}
En el paso 2 de nuestro código, usamos la segunda opción:
\begin{equation}
    \dot{R} = R \begin{bmatrix} 0 & -\omega_3 & \omega_2 \\ \omega_3 & 0 & -\omega_1 \\ -\omega_2 & \omega_1 & 0 \end{bmatrix}
\end{equation}
Esto es extremadamente conveniente porque las componentes $\omega_1, \omega_2, \omega_3$ son precisamente las que ya calculamos en el Paso 1 al resolver las ecuaciones de Euler. 

Integrar esta ecuación diferencial nos permite "arrastrar" la orientación del cuerpo paso a paso, asegurando que los ejes $\hat{e}_i$ siempre sean ortogonales y representen fielmente la rotación física del objeto, permitiéndonos luego calcular $\vec{\omega}_{fijo}$ simplemente multiplicando $R$ por $\vec{\omega}_{cuerpo}$.

\end{document}
